\documentclass[11pt]{article}
\usepackage{amsmath,amssymb,mathtools}
\usepackage{fontspec,unicode-math}
\setmainfont{Latin Modern Roman}
\setsansfont{Latin Modern Sans}
\setmathfont{Latin Modern Math}
\usepackage{xurl,hyperref}
\newcommand{\sref}[2]{\hyperlink{source:#1:#2}{[S#2]}}
\newcommand{\cataloguescope}[1]{\paragraph{Mathematical question.}#1}
\begin{document}
\section{Yao-Yao Graph a Spanner?}\label{statement}

\subsection{Definitions and mathematical statement}

\paragraph{Shared notation.}$\mathbb N=\{1,2,\ldots\}$, and $\mathbb Z,\mathbb Q,\mathbb R,\mathbb C$ denote the integers, rationals, reals and complex numbers. Any other conventions are specified locally.


Let $S\subset\mathbb R^2$ be finite and fix $k$ equally spaced ray directions. Translate the resulting $k$ cones to each $u\in S$. The directed Yao graph $Y_k$ selects, in each nonempty cone at $u$, a shortest outgoing edge $u\to v$. To form $YY_k$, at each vertex and in each cone retain only the shortest incoming edge from that cone, then forget the edge orientations. Use a consistent convention for boundary rays and equal distances. A geometric graph is a $t$-spanner if its shortest Euclidean-edge-length path distance satisfies $d_G(u,v)\le t\lVert u-v\rVert$ for every pair.
\paragraph{Statement recorded in the supplied source.}
For which fixed cone counts $k$ is the Yao--Yao graph $YY_k$ a $t$-spanner with a constant $t$ independent of the finite point set? This records the source\textquotesingle s question for the family $YY_k$, allowing the constant to depend on $k$. \sref{statement}{2}
\subsection{Short English statement}
For which fixed cone counts $k$ is the Yao--Yao graph $YY_k$ a $t$-spanner with a constant $t$ independent of the finite point set? This records the source\textquotesingle s question for the family $YY_k$, allowing the constant to depend on $k$.
\subsection{Sources}
\begin{description}
\item[\hypertarget{source:statement:1}{S1}] ULAM.AI and the UnsolvedMath Contributors, UnsolvedMath: A Curated Collection of Open Mathematics Problems, catalogue entry. CC BY 4.0. \url{https://huggingface.co/datasets/ulamai/UnsolvedMath}.
\item[\hypertarget{source:statement:2}{S2}] Erik D. Demaine, Joseph S. B. Mitchell and Joseph O\textquotesingle Rourke (editors), \emph{The Open Problems Project}, maintained original source, Problem 70, Statement and complete Partial/Related Results. \url{https://github.com/edemaine/topp/blob/main/Problems/P.000070}.
\end{description}
\label{end:statement}
\end{document}
